\subsection{Resolutions}

In what follows, one always identifies a module with the complex of which it is the degree-zero component and of which all other components are zero.

DEFINITION 1. --- Let M be an A-module. A left resolution of M is a pair (P, p) where P is a complex zero to the right and p : P → M is a homologism. A right resolution of M is a pair (e, E) where E is a complex zero to the left and e : M → E a homologism.

One calls the length of the resolution (P, p) (resp. (e, E)) the length of the complex P (resp. E). If (P, p) and (P', p') (resp. (e, E) and (e', E')) are two left (resp. right) resolutions of M, a morphism of complexes \(f: P \to P'\) such that \(p' \circ f = p\) (resp. \(g: E \to E'\) such that \(g \circ e = e'\) ) is called a morphism of resolutions.

PROPOSITION 2. --- Let P be a complex that is zero on the right and let p : P → M be a morphism. For (P, p) to be a left resolution of M, it is necessary and sufficient that the sequence

\[
\dots \longrightarrow \mathrm{P} _ {n} \xrightarrow {d _ {\mathrm{P}}} \mathrm{P} _ {n - 1} \longrightarrow \dots \quad \longrightarrow \mathrm{P} _ {1} \xrightarrow {d _ {\mathrm{P}}} \mathrm{P} _ {0} \xrightarrow {p _ {0}} \mathrm{M} \longrightarrow 0\tag{1}
\]

be exact.

Indeed, to say that \(p: \mathbf{P} \to \mathbf{M}\) is a homologism means that \(\mathrm{H}_{i}(\mathbf{P}) = 0\) for \(i > 0\) and that \(p_{0}\) induces an isomorphism from Coker ( \(d_{\mathbf{P}}: \mathbf{P}_{1} \to \mathbf{P}_{0}\) ) onto \(\mathbf{M}\) .

Similarly:

PROPOSITION 2bis. --- Let E be a complex that is zero on the left and let e : M → E be a morphism. For (e, E) to be a right resolution of M, it is necessary and sufficient that the sequence

\[
0 \longrightarrow \mathrm{M} \stackrel {e _ {0}} {\longrightarrow} \mathrm{E} ^ {0} \stackrel {d _ {\mathrm{E}}} {\longrightarrow} \mathrm{E} ^ {1} \longrightarrow \dots \longrightarrow \mathrm{E} ^ {n} \stackrel {d _ {\mathrm{E}}} {\longrightarrow} \mathrm{E} ^ {n + 1} \longrightarrow \dots\tag{1 bis}
\]

be exact.

By abuse of language, one often says that the sequence (1) (resp. (1 bis)) is a left (resp. right) resolution of M.

DEFINITION 2. --- A projective (resp. free, resp. flat) resolution of the A-module M is a left resolution (P, p) of M such that the complex P is projective (resp. free, resp. flat) (X, p.~25). An injective resolution of M is a right resolution (e, E) of M such that the complex E is injective (loc. cit.).

Examples. --- 1) Suppose that the ring A is principal; let M be an A-module and \((x_{i})_{i\in I}\) a generating family of M. Let \(L_{0}\) be the free module \(\mathbf{A}^{(I)}\) , \((e_{i})\) its canonical basis, and define \(p:L_{0}\to\mathbf{M}\) by \(p(e_{i})=x_{i}\) . The morphism \(p\) is surjective and its kernel \(L_{1}\) is a free A-module by VII, § 3, cor. 2 to th. 1, so the exact sequence

\[
0 \to \mathrm{L} _ {1} \to \mathrm{L} _ {0} \xrightarrow {p} \mathrm{M} \to 0
\]

is a free resolution of M of length 1. If I is finite, \(L_{0}\) and \(L_{1}\) are finitely generated.

\begin{enumerate}
\def\labelenumi{\arabic{enumi})}
\setcounter{enumi}{1}
\tightlist
\item
  Suppose A is commutative; let E be an A-module and u an endomorphism of E. Denote by \(E_{u}\) the A{[}X{]}-module obtained by endowing E with the structure defined by
\end{enumerate}

\[
(p, x) \mapsto p (u) (x) \quad \text { for } \quad p \in \mathrm{A} [ \mathrm{X} ] \quad \text { and } \quad x \in \mathrm{E}.
\]

By III, p.~106, we have an exact sequence:

\[
0 \to \mathrm{A} [ \mathrm{X} ] \otimes_ {\mathrm{A}} \mathrm{E} \stackrel {\psi} {\to} \mathrm{A} [ \mathrm{X} ] \otimes_ {\mathrm{A}} \mathrm{E} \stackrel {\varphi} {\to} \mathrm{E} _ {u} \to 0
\]

where \(\varphi(p \otimes x) = p.x\) and \(\psi(p \otimes x) = Xp \otimes x - p \otimes u(x)\) for \(p \in A[X]\) and \(x \in E\) . This exact sequence is a resolution of length 1 of \(E_u\) free (resp. projective, resp. finitely generated) if E is a free (resp. projective, resp. finitely generated) A-module. 3) If A is principal, the exact sequence

\[
0 \rightarrow \mathrm{A} \rightarrow \mathrm{K} \rightarrow \mathrm{K/A} \rightarrow 0
\]

is an injective resolution of length 1 of the A-module \(\mathbf{A}_s\) (X, p.~18, example 1).

PROPOSITION 3. — Let \(f: \mathbf{M}' \to \mathbf{M}\) be a homomorphism of A-modules, \(p': \mathbf{P}' \to \mathbf{M}'\) be a morphism into \(\mathbf{M}'\) of a complex that is zero on the right and projective \(\mathbf{P}'\) , and \(p: \mathbf{P} \to \mathbf{M}\) a left resolution of M. There exists a morphism of complexes \(\tilde{f}: \mathbf{P}' \to \mathbf{P}\) and one only up to homotopy, such that \(p \circ \tilde{f} = f \circ p'\) .

Consider the complex \(\overline{\mathbf{P}}\) defined as follows: \(\overline{\mathbf{P}}_n = \mathbf{P}_n\) for \(n \neq -1\) , \(\overline{\mathbf{P}}_{-1} = \mathbf{M}\) , \(d_{\overline{\mathbf{P}},n} = d_{\mathbf{P},n}\) for \(n \neq 0, -1\) , \(d_{\overline{\mathbf{P}},0} = p_0\) , \(d_{\overline{\mathbf{P}},-1} = 0\) , and the complex \(\overline{\mathbf{P}}'\) defined analogously. Applying to the complexes \(\overline{\mathbf{P}}\) and \(\overline{\mathbf{P}}'\) prop. 1 a) with \(r=0\) , \(u_i = 0\) for \(i < -1\) and \(u_{-1} = f\) , we obtain prop. 3.

COROLLARY. — Let (P, p) and (P', p') be two projective resolutions of M. There exists one and only one homotopism up to homotopy \(\alpha: P' \to P\) such that \(p \circ \alpha = p'\) .

Indeed, there exists a morphism \(\alpha: P' \to P\) (resp. \(\beta: P \to P'\) ) such that \(p \circ \alpha = p'\) (resp. \(p' \circ \beta = p\) ). Since \(p \circ \alpha \circ \beta = p\) (resp. \(p' \circ \beta \circ \alpha = p'\) ), \(\alpha \circ \beta\) is homotopic to \(1_P\) (resp. \(\beta \circ \alpha\) is homotopic to \(1_{P'}\) ).

PROPOSITION 3 bis. — Let \(g: \mathbf{N} \to \mathbf{N}'\) be a homomorphism of A-modules, \(e': \mathbf{N}' \to \mathbf{E}'\) a morphism of \(\mathbf{N}'\) into a complex that is zero on the left and injective \(\mathbf{E}'\) , and \(e: \mathbf{N} \to \mathbf{E}\) a right resolution of \(\mathbf{N}\) . There exists a morphism of complexes \(\tilde{g}: \mathbf{E} \to \mathbf{E}'\) and one only up to homotopy, such that \(\tilde{g} \circ e = e' \circ g\) .

This is proved like Prop. 3 with the aid of Prop. 1 b).

COROLLARY. --- Let \((e, E)\) and \((e', E')\) be two injective resolutions of N; there exists a homotopism \(\alpha : E \to E'\) and only one up to homotopy such that \(\alpha \circ e = e'\) .

